\section{Considerations for Mixed Effects Models}

Because the dataset in question includes categorical variables for governorate (\emph{adm1}) and district (\emph{adm2}), with 22 and 231 levels (resp.), a fixed effects model that uses all variables is out of the question. As a result, it would be desirable to generalize the approach previously detailed to mixed effects models, in which categorical variables with many levels can be treated as random effects.

The model to be fit is now expressed as follows:

\begin{equation}
    \label{eq:mixed_effects_ts}
    \yy_{ti} = x_f\beta + x_r\mathbf{b} + x_t\uu + \varepsilon, \qquad \varepsilon \sim N(\mathbf{0},\sigma_\varepsilon^2I)
\end{equation}

In this model the vector $\mathbf{b}$ corresponds to the random effects of each governorate and district. 

Let's briefly summarize how mixed effects models are fit to establish any ramifications for integrating such a model along with a time series into our EM algorithm.